\setlength{\baselineskip}{20pt}
\begin{abstract}
随着自动泊车与智能驾驶技术的快速发展，停车位检测作为环境感知系统的核心环节，其检测精度与鲁棒性直接决定了自动泊车系统的安全性与可靠性。与传统目标检测任务不同，停车位检测具有强结构约束特性，同时面临标注成本高昂、场景分布呈现长尾效应等挑战，导致其在复杂真实环境下仍存在数据规模受限、模型泛化能力不足以及检测结果结构不一致等问题。针对上述挑战，本文围绕停车位检测任务，从数据构建与学习范式、端到端结构建模以及多帧时序信息融合三个层面开展研究:

（1）针对现有数据集场景覆盖有限、人工标注成本高昂的瓶颈，本文提出了一种面向复杂场景停车位检测的半监督学习方法。为了更好的建模真实世界中的停车位检测场景，
本文首先构建了包含 35739 张图像的真实场景停车位检测数据集 CRPS-D，其具备场景类型丰富、停车位密度高以及异常车位比例大等特点。在此基础上，结合停车位关键点分布稀疏且结构约束显著的任务特性，本文提出了一种基于教师–学生架构的停车位检测半监督学习方法（SS-PSD）。SS-PSD 通过置信度引导掩码一致性与自适应特征扰动机制，有效挖掘了无标注数据中的空间结构先验。实验结果表明，在仅使用 1/12 标注训练数据的情况下，模型在 ps2.0 数据集上的检测平均精度（AP）达到 90.47\%，在 CRPS-D 数据集上的 AP 达到 79.75\%，在相同标注比例设置下显著优于对应的全监督基线模型。

（2）针对现有方法依赖人工设计后处理规则、难以保证结构一致性的局限，本文提出了一种面向结构感知的停车位关键点检测与匹配方法（SAKM-PSD）。SAKM-PSD 在网络内部显式建模几何拓扑约束，实现了停车位关键点定位与结构关系的联合优化，从而摒弃了传统启发式匹配流程。实验结果表明，SAKM-PSD 在 ps2.0 数据集上取得了 99.37\% 的查准率、99.47\% 的查全率以及 99.47\% 的 AP；在 CRPS-D 数据集上，取得了 92.21\% 的查准率、85.54\% 的查全率以及 83.58\% 的 AP。

（3）针对单帧检测易受瞬时遮挡和环境噪声干扰的问题，本文提出了一种面向多帧时序信息融合的停车位检测方法（MTF-PSD），将静态空间建模扩展至时空联合建模。MTF-PSD 通过挖掘连续视频序列中的时序相关性，对结构感知特征进行递归建模与时序平滑，有效提升了检测结果在时间维度上的一致性。实验结果表明，MTF-PSD 在 ps2.0 数据集上取得了 99.42\% 的查准率、99.37\% 的查全率以及 98.64\% 的 AP；在自建时序数据集 PSTD 上，取得了 92.53\% 的查准率、92.33\% 的查全率以及 91.97\% 的 AP。


	\noindent\keywords{停车位检测；数据集；半监督学习；关键点检测；结构感知；时序建模}
\end{abstract}